\chapter{Наследование второго семейного тензора}
\label{ch:v8-second-family-tensor-inheritance-no-go-gate}

\section{Три уже вычисленных кандидата}

База проекта действительно содержит операторы, чувствительные к семейным
осям. Но для открытия нового тома недостаточно существования матрицы: нужен
канонически выбранный оператор внутри текущего родителя и слепой успешный
тест после заморозки выбора.

Первый кандидат --- общий $J$-парный прямоугольник Краевского. Он выводит
\begin{equation}
 \Tr D_F^4=104+16\cos\theta\,\Tr(P_-H_uH_d),
\end{equation}
но предполагает ненулевой общий коннектор. Ветка такого коннектора уже
закрыта подготовительными гейтами, а два минимума остаются CP-чётными и
ориентационно вырожденными.

Второй кандидат --- моментный коммутатор
\begin{equation}
 K_s=\frac{[Z_{4,s},Z_{6,s}]}{2i}.
 \label{eq:v8-moment-commutator-candidate}
\end{equation}
Это настоящий независимый чувствительный к собственным векторам тензор. Однако ранее
замороженная по массам ветвь не прошла даже массового ряда, а в слепом CKM-тесте
дала $s_{12}>0.99$. Повторный выбор знака или размещения после просмотра
смешивания запрещён.

Третий кандидат --- меню аффинных операторов инцидентности. Двенадцать из
шестнадцати направлений порождают $M_3(\mathbb C)$, но геометрия не выбирает
ни одно из них. Усреднение по остаточной симметрии возвращает редуцируемую
алгебру размерности пять.

\section{Сопоставление с Томом VII}

Единственная голономия Тома VII задаёт классовую функцию: она видит
собственные фазы, но не собственные оси. Поэтому ни один из трёх старых
кандидатов не наследуется автоматически текущим качественным родителем:
первому нужен закрытый коннектор, второй уже провалил замороженный слепой
тест, третий остаётся невыбранным меню.

\begin{theorem}[Запрет наследования второго семейного тензора]
В записанной базе существуют независимые семейные операторы, но нет ни
одного, который одновременно выводится текущей архитектурой, канонически
выбирается до просмотра CKM/PMNS и имеет прошедшую слепую ветвь. Поэтому
примитив B нельзя считать уже найденным или без изменений перенести в
следующий том.
\end{theorem}

\begin{remark}
Это не универсальный запрет новой геометрии. Он закрывает только повторное
использование уже проверенного меню как будто оно является новым выводом.
\end{remark}

\begin{protocol}
\begin{enumerate}
 \item межсекторный цикл: \textbf{оператор есть, коннектор предположен};
 \item коммутатор момента: \textbf{независим, слепой тест провален};
 \item набор инцидентностных операторов: \textbf{12 рабочих операторов, селектора нет};
 \item одна голономия Тома VII выбирает семейные оси: \textbf{нет};
 \item наследуемый примитив B: \textbf{не получен};
 \item примитивы A, B, C текущей архитектуры: \textbf{закрыты};
 \item Том VIII на старой базе: \textbf{не открывать};
 \item допустимое продолжение: \textbf{только новый типизированный оператор
 с заранее объявленным стоп-тестом}.
\end{enumerate}
\end{protocol}

Машинный сертификат сохранён в
\path{s2t_v8_second_family_tensor_inheritance_no_go_gate_results.json}.